Un tub de buit com el de la figura adjunta té l'ànode A fet de coure i la distància entre els elèctrodes és $d = 30\ \mathrm{cm}$. Establim un camp elèctric uniforme de A a B que genera una diferència de potencial de 3~V i il·luminem l'ànode amb radiacions que tenen fotons incidents amb una energia de 10~eV. Observem que al càtode B arriben electrons amb una energia cinètica de 2,3~eV.

\figura[0.4]{fig1.png}

\begin{apartat}{1}
Quina és la freqüència i la longitud d'ona de la radiació incident (expressada en nm)? Quin és el valor del camp elèctric $E$?
\end{apartat}

\begin{apartat}{1}
Amb quina energia cinètica surten emesos els electrons arrencats de l'ànode A? Quin és el treball d'extracció del coure en eV?
\end{apartat}

\dades{%
  $h = 6,626\times 10^{-34}\ \mathrm{J\,s}$\\
  $Q_\text{electró} = -1,602\times 10^{-19}\ \mathrm{C}$\\
  $c = 3\times 10^{8}\ \mathrm{m\,s^{-1}}$\\
  $1\ \mathrm{eV} = 1,602\times 10^{-19}\ \mathrm{J}$}
